\section{术语与指标速查}

\subsection{关键术语}

\begin{figure}[H]
\centering
\includegraphics[width=0.65\textwidth]{slides-images/slide-39.jpg}
\caption{术语与指标速查 slide（slide-39）。}
\end{figure}

\begin{table}[H]
\centering
\begin{tabular}{p{4cm}p{7cm}}
\toprule
术语 & 释义 \\
\midrule
CoT & Chain-of-Thought，向模型展示完整推理步骤 \\
Self-consistency & 多路径采样后取众数，减少偶发错误 \\
Governance loop & Prompt log + dashboard + incident grade-book \\
KV cache & Key/value 缓存，用于 decoder-only 多轮对话加速 \\
\bottomrule
\end{tabular}
\caption{Lecture 15 重复出现的关键术语。}
\end{table}

\begin{knowledgebox}{术语记忆建议}
把关键术语写进 `glossary.md` 并在 incident retro 时复盘，帮助新加入的工程师快速上手。
\end{knowledgebox}

\subsection{指标矩阵与 dashboard}

\begin{table}[H]
\centering
\begin{tabular}{p{4cm}p{5cm}}
\toprule
指标 & 作用 \\
\midrule
Self-consistency score & 衡量 prompt/plan 稳定性 \\
Hallucination rate & 记录 fabrication 事件占比 \\
Token budget & 每次请求的 token/compute 限额 \\
Attention entropy & 判断模型是否跳到不同模式 \\
\bottomrule
\end{tabular}
\caption{治理与调试时常用的指标。}
\end{table}

\begin{warningbox}{指标滞后带来的风险}
\begin{itemize}
    \item 过度依赖 accuracy 可能掩盖 hallucination；
    \item Hallucination rate 如果 label delay，会让 governance 反应滞后；
    \item 没有 attention entropy 监控，很难确认 emergent jump 是否真实。
\end{itemize}
\end{warningbox}

\subsection{本章小结}

把术语与指标写下来，便于新成员快速理解讲座用语，也可以在 incident review 中快速定位 key term + metric。

